\section{Related Work}

\paragraph{Virtual unwrapping of the Herculaneum scrolls.} Virtual unwrapping was introduced by Seales et al.~\shortcite{seales2016damage}, who recovered text from the En-Gedi scroll by segmenting its surface within a CT volume and flattening it to a plane. Applying the same idea to the far more tightly-wound Herculaneum scrolls has driven several recent automatic approaches, collected by the Vesuvius Challenge. {\em Spiral-fitting} \cite{henderson2026virtually} uses the fact that a scroll is, in its original form, a single rolled sheet, and fits an extruded spiral to the prediction by solving a variational problem; this is elegant, but degrades wherever the neural-network prediction is too noisy for any single spiral to fit. {\em Surface-tracing} approaches find small patches on the surface prediction and grow them iteratively, turning the problem into an advancing front. To the best of our knowledge, the open difficulty with this approach is how adjacent patches are merged and how the resulting topological artifacts are resolved, especially when the merge is performed after flattening the patches to two dimensions. {\em ThaumatoAnakalyptor} extracts a point cloud using blurring and Sobel edge detection, classifies the points into wraps, and reconstructs the surface by solving a winding-angle assignment problem, sharing the prediction-noise sensitivity of spiral-fitting. ScrollFiesta also grows an advancing front via ball-pivoting, but handles topology directly on the reconstructed mesh rather than in raw prediction space. In addition to avoiding the necessity of complicated front merges, this in-situ approach makes topological repair such as hole filling and bridge splitting feasible. When we do assign a winding index for the final unwrapping (Section~\ref{sec:unwrap}), we adopt the diffeomorphic-spiral idea of these methods, but apply it to the {\em repaired} mesh in a slice domain where the residual fusion bridges can be cut geometrically, rather than to the noisy prediction directly.

\paragraph{Surface reconstruction and meshing.} Surface reconstruction is a mature field, but its standard tools are a poor fit here because they are built to produce {\em two-sided, watertight} surfaces. Marching cubes~\cite{lorensen1998marching} and its descendants, such as dual contouring~\cite{ju2002dual} and the differentiable FlexiCubes~\cite{shen2023flexicubes}, extract a closed isosurface; Poisson reconstruction~\cite{kazhdan2006poisson} fits a closed indicator function and, in its distributed form~\cite{kazhdan2023distributed}, scales to large volumes. The main problem with these classical approaches is that a scroll sheet is an open, single-sided manifold: applying methods designed for closed surfaces yields two-sided envelopes that must subsequently be converted to single-sided sheets. Before designing our current approach, we investigated the possibility of meshing the Herculaneum papyri using both marching cubes and Poisson surface reconstruction as starting points. These failures motivated our current algorithm.

\paragraph{Mesh parameterization and atlas construction.} Flattening a surface to the plane is a mature area, but the classical objectives are the wrong ones here. Angle-based flattening and LSCM~\cite{sheffer2005abf++} minimize angular or conformal distortion and offer {\em no} guarantee that the resulting map is injective; on a long, high-aspect, physically damaged strip, folded and overlapping triangles are a routine outcome rather than a corner case. That is exactly the guarantee a scroll needs, so where ScrollFiesta hands a clean disk-like patch to an external flattener it uses a symmetric-Dirichlet energy~\cite{smith2015bijective}, which is locally injective by construction, and it never uses a conformal map as the final authority. A surface that will not flatten as one piece is instead cut into charts and packed into an atlas, and a body of work addresses making that packing compact and low-distortion~\cite{limper2018box}.

What all of these share is an assumption that the input surface is correct and that the remaining problem is distortion. A scroll violates that assumption twice, and both violations are structural rather than numerical. First, the correct layout is not the one that minimizes distortion. It is the one in which a chart sits exactly one circumference away from the chart one wrap inside it. That constraint is global, integer-valued, and external to the mesh: no local energy on the surface can express it, and a distortion-optimal flattening of a wrongly fused pair of wraps is still wrong. Second, the failure to avoid is not per-chart injectivity but {\em inter-chart collision} --- two different wraps occupying the same $(u,v)$ --- which a free-boundary packer has no reason to prevent, since in a conventional atlas a chart may be placed anywhere that is empty. Our slice-domain construction attacks both directly: developed length comes from measured arc length along cut strokes and is made monotone by isotonic projection, so a chart is overlap-free in $u$ by construction rather than by penalty, and the remaining inter-chart collisions are then audited exactly rather than approximated by an energy. The closest existing formulation is StrokeStrip~\cite{pagurek2021strokestrip}, which fits one common parameterization to a {\em cluster} of strokes rather than to a single connected patch; we adopt its slice-domain view and add the integer winding assignment a scroll additionally requires, solving the two as alternating discrete and continuous stages (Section~\ref{sec:unwrap}).

